% 图契约：相邻节点公式：先叠加支路电流，再按端口归一化；图中的端口、几何、箭头和方程对应本题。
% 版面：沿用前三模块的几何图+推导分区；目标印宽17cm，简图至中等密度。
% 布局：左侧真实结构，右侧关键关系；真实电路边不作为架构流程箭头处理。
\documentclass[tikz,border=7pt]{standalone}
\usepackage{fontspec,xeCJK,amsmath,amssymb}
\IfFontExistsTF{Times New Roman}{\setmainfont{Times New Roman}}{\setmainfont{TeX Gyre Termes}}
\IfFontExistsTF{Heiti SC}{\setCJKmainfont[AutoFakeBold=1.5]{Heiti SC}}{\setCJKmainfont[AutoFakeBold=1.5]{FandolHei-Regular}}
\usetikzlibrary{arrows.meta,bending,calc,patterns,positioning,decorations.pathmorphing}
\definecolor{ink}{HTML}{243745}
\definecolor{blue}{HTML}{275DAB}
\definecolor{green}{HTML}{177C63}
\definecolor{red}{HTML}{BD4935}
\definecolor{orange}{HTML}{AA6C16}
\definecolor{purple}{HTML}{715393}
\tikzset{every node/.style={font=\fontsize{9.3}{12.5}\selectfont,text=ink},
arrow/.style={-{Stealth[length=5pt 1.5,width'=0pt 3]},shorten >=2pt,shorten <=1pt,line width=1pt},
arrow short/.style={-{Stealth[length=3pt 1.0,width'=0pt 2.5]},shorten >=1pt,shorten <=0pt,line width=.8pt},
axis/.style={arrow,draw=ink!60,line width=.6pt},
vec/.style={arrow,draw=blue},
wire/.style={draw=ink!75,line width=.8pt},
guide/.style={dashed,draw=ink!45,line width=.5pt},
card/.style={rounded corners=2pt,draw=blue!65,fill=blue!4,align=center,text width=6.5cm,inner xsep=9pt,inner ysep=7pt},
title/.style={font=\fontsize{10.5}{14}\selectfont\bfseries,text width=15cm,align=center},
label/.style={font=\fontsize{8.4}{11}\selectfont,align=center},
note/.style={label,text width=6.1cm}}
\begin{document}
\begin{tikzpicture}[x=1cm,y=1cm]
\path[use as bounding box] (-8,-3.5) rectangle(8,3.65);
\node[title] at (0,3.25) {\quad 相邻节点公式：先叠加支路电流，再按端口归一化};
\begin{scope}[shift={(-5.2,.6)}]\draw[ink,fill=white](0,0)circle(.3);\node at(0,0){$A$};\draw[vec,green](0.38,0.0)--(1.4,0.0);\draw[vec,green](0.19000000000000006,0.32908965343808666)--(0.7000000000000001,1.212435565298214);\draw[vec,green](-0.18999999999999995,0.3290896534380867)--(-0.6999999999999997,1.2124355652982142);\draw[vec,green](-0.38,4.6536578367599424e-17)--(-1.4,1.7145055188062944e-16);\draw[vec,green](-0.19000000000000017,-0.3290896534380866)--(-0.7000000000000006,-1.2124355652982137);\draw[vec,green](0.19000000000000006,-0.32908965343808666)--(0.7000000000000001,-1.212435565298214);\node[label]at(0,1.9){A 点平衡散流};\end{scope}\node[card,text width=4.1cm,draw=green!65,fill=green!4]at(-5.2,-1.35){$I_{AB}^{(1)}=I/d$};
\begin{scope}[shift={(0,.6)}]\draw[ink,fill=white](0,0)circle(.3);\node at(0,0){$B$};\draw[vec,blue](1.4,0.0)--(0.38,0.0);\draw[vec,blue](0.7000000000000001,1.212435565298214)--(0.19000000000000006,0.32908965343808666);\draw[vec,blue](-0.6999999999999997,1.2124355652982142)--(-0.18999999999999995,0.3290896534380867);\draw[vec,blue](-1.4,1.7145055188062944e-16)--(-0.38,4.6536578367599424e-17);\draw[vec,blue](-0.7000000000000006,-1.2124355652982137)--(-0.19000000000000017,-0.3290896534380866);\draw[vec,blue](0.7000000000000001,-1.212435565298214)--(0.19000000000000006,-0.32908965343808666);\node[label]at(0,1.9){B 点平衡汇流};\end{scope}\node[card,text width=4.1cm,draw=blue!65,fill=blue!4]at(0,-1.35){$I_{AB}^{(2)}=I/d$};
\begin{scope}[shift={(5.2,.6)}]\draw[wire,draw=ink](-1.6,0)--(-0.275,0.0) (0.275,0.0)--(1.6,0);\begin{scope}[shift={(0.0,0.0)},rotate=0.0]\draw[draw=ink,fill=white,line width=.8pt](-0.275,-.085)rectangle(0.275,.085);\end{scope}
\node[label]at(0.0,0.42){$r$};
\fill[ink] (-1.6,0)circle(2pt);\node[below=4pt] at(-1.6,0){$A$};
\fill[ink] (1.6,0)circle(2pt);\node[below=4pt] at(1.6,0){$B$};
\draw[vec,red](-1.2,.65)--(1.2,.65)node[midway,above]{$2I/d$};\node[label]at(0,1.9){叠加后的直接支路};\end{scope}\node[card,text width=4.1cm,draw=red!65,fill=red!4]at(5.2,-1.35){$U_{AB}=2Ir/d$\\$I_{\text{端}}=NI/(N-1)$};
\node[card,text width=14.3cm,draw=blue!65,fill=blue!4]at(0,-2.65){$R_{AB}=\dfrac{2(N-1)}{Nd}r$；前提：相邻d条支路在本辅助态下等价。放射支路仅示意一般d。};

\end{tikzpicture}\end{document}